% Generated from the audited Markdown source; responsible author: Carptopus.
\documentclass[11pt]{article}
\usepackage[a4paper,margin=27mm]{geometry}
\usepackage[T1]{fontenc}
\usepackage{lmodern}
\usepackage{microtype}
\usepackage{amsmath,amssymb,amsthm,mathtools}
\usepackage{needspace}
\usepackage{xcolor}
\usepackage[colorlinks=true,linkcolor=blue!55!black,citecolor=blue!55!black,urlcolor=blue!65!black]{hyperref}
\usepackage{fancyhdr}

\newtheorem{theorem}{Theorem}[section]
\newtheorem{lemma}[theorem]{Lemma}
\newtheorem{proposition}[theorem]{Proposition}
\newtheorem{corollary}[theorem]{Corollary}
\numberwithin{equation}{section}

\pagestyle{fancy}
\fancyhf{}
\fancyhead[L]{Carptopus}
\fancyhead[R]{Palindromicity of antichain-box polynomials}
\fancyfoot[C]{\thepage}
\setlength{\headheight}{14pt}
\setlength{\parindent}{0pt}
\setlength{\parskip}{0.45em}
\setlength{\emergencystretch}{4em}
\urlstyle{same}

\title{Palindromicity of antichain polynomials\\
of three-dimensional boxes}
\author{Carptopus\\\href{mailto:carptopus@163.com}{carptopus@163.com}}
\date{29 August 2026}

\hypersetup{
  pdftitle={Palindromicity of antichain polynomials of three-dimensional boxes},
  pdfauthor={Carptopus},
  pdfsubject={Complete palindromicity classification for antichain generating polynomials of three-dimensional box posets},
  pdfkeywords={antichain generating polynomial, palindromicity, gamma-positivity, minuscule poset, last passage, lattice path, vertex separator}
}

\begin{document}
\maketitle

\begin{center}
\fcolorbox{black!35}{black!4}{\parbox{0.88\textwidth}{\centering\small
Internal candidate manuscript, version 0.1-beta. The mathematical proof and the
complete Markdown manuscript have passed project-internal zero-trust audits.
External mathematical review and formal peer review are pending.}}
\end{center}

\begin{abstract}


For positive integers $a,b,c$, let

\nopagebreak[4]
\[
N_{a,b,c}(x)=\sum_A x^{|A|}
\]

be the antichain generating polynomial of the product poset
$[a]\times[b]\times[c]$. We prove that, after sorting
$a\le b\le c$,

\nopagebreak[4]
\[
N_{a,b,c}(x)\ \text{is palindromic}
\quad\Longleftrightarrow\quad
(a,b,c)=(1,r,r)\ \text{or}\ (2,r,r+1).
\]

The new part is the exclusion of all boxes with $a,b,c\ge3$. We use the
last-passage representation of antichains by nonnegative matrices. In the
monic parameter range

\nopagebreak[4]
\[
c=a+b-1-2d,
\]

we classify all binary supports one below maximum and all possible weight
increments. For $d\ge1$ this gives the exact coefficient

\nopagebreak[4]
\[
[x^{w-1}]N_{a,b,c}(x)
=ab+2+(a-d)(a-d-1)+(b-d)(b-d-1),
\]

where $w=ab-d(d+1)$. This coefficient is strictly smaller than
$[x]N_{a,b,c}(x)=abc$. The remaining case $d=0$ has an explicit coefficient
strictly larger than $abc$. Thus the first and next-to-leading coefficients
exclude palindromicity in every genuine higher-layer case. Combined with the
known two-layer classification, this proves the rectangular minuscule branch
of a conjecture of Ding and Dong.

\end{abstract}

\section{Introduction}


For a finite poset $P$, its antichain generating polynomial is

\nopagebreak[4]
\[
N_P(x)=\sum_{A\in\mathcal A(P)}x^{|A|},
\]

where $\mathcal A(P)$ denotes the set of antichains of $P$. Ding and Dong
studied palindromicity, real-rootedness, and gamma-positivity of these
polynomials for minuscule posets \cite{ding2019}. Their Conjecture B states that, for a
connected minuscule poset $P$, palindromicity of
$N_{[k]\times P}(x)$ implies strict gamma-positivity.

The rectangular minuscule family takes

\nopagebreak[4]
\[
P=[b]\times[c],
\]

so the relevant polynomial is that of the three-dimensional box
$[a]\times[b]\times[c]$. The one-layer case is elementary. The two-layer
case was recently classified together with a proof of real-rootedness and
strict gamma-positivity \cite{carptopus2026}. The purpose of this paper is to settle every
remaining number of layers.

Because the three factors are symmetric, we always sort

\nopagebreak[4]
\[
a\le b\le c.
\]

Our main result is the following.

\Needspace{0.25\textheight}
\begin{theorem}

For positive integers $a\le b\le c$, the polynomial $N_{a,b,c}(x)$ is
palindromic if and only if

\nopagebreak[4]
\[
(a,b,c)=(1,r,r)
\quad\text{or}\quad
(a,b,c)=(2,r,r+1).
\]

\end{theorem}

The genuinely new assertion is:

\Needspace{0.25\textheight}
\begin{theorem}

If $3\le a\le b\le c$, then $N_{a,b,c}(x)$ is not palindromic.

\end{theorem}

The proof needs only two coefficients, but obtaining the second one requires
a structural classification. Amanov and Yeliussizov \cite{amanov2023} give a last-passage
correspondence between plane partitions and nonnegative matrices. Under this
correspondence, the size of an antichain becomes the number of positive
matrix entries. A near-maximum antichain is therefore a near-maximum support
subject to a last-passage constraint.

The complement of the maximum support consists of two triangular end regions.
Removing points from these triangles creates path deficits. We quantify the
minimum cost of compensating those deficits inside the central band. A
frontier recurrence first forces all removals to form one interval on the
outermost diagonal of one end. A minimum-separator theorem then shows that the
only compensation is the adjacent shadow, apart from two full cross-end
exchanges. This yields exactly $ab+2$ binary near-maximum supports. A separate
dominator argument classifies all supports that can carry one additional
unit of weight.

Section 2 recalls the last-passage model and the monic parameter reduction.
Sections 3--5 prove the path-compensation classification. Section 6 classifies
weight increments and obtains the next-to-leading coefficient for $d\ge1$.
Section 7 treats $d=0$. Section 8 proves the complete classification.

\section{Last passage and the monic parameter range}


Write

\nopagebreak[4]
\[
Q=[0,a-1]\times[0,b-1].
\]

For a nonnegative integer matrix $M=(M_{ij})_{(i,j)\in Q}$, let $L(M)$ be
the maximum of

\nopagebreak[4]
\[
\sum_{(i,j)\in P}M_{ij}
\]

over all northeast lattice paths $P$ from $(0,0)$ to $(a-1,b-1)$.

Amanov and Yeliussizov [2, Theorem 3.1, equation (3), Corollary 3.2(i),
Lemma 3.4, and Lemma 4.2] give a bijection under which

\nopagebreak[4]
\[
N_{a,b,c}(x)
=\sum_{\substack{M\in\mathbb Z_{\ge0}^{a\times b}\\L(M)\le c}}
x^{|\operatorname{supp}M|}.
\tag{2.1}
\]

Indeed, their inverse map satisfies

\nopagebreak[4]
\[
M_{ij}>0
\quad\Longleftrightarrow\quad
(i,j,\pi_{ij})
\text{ is a top corner of the corresponding plane partition}.
\]

The top corners are the maximal elements of the corresponding order ideal,
and hence form its antichain of maximal elements. This proves (2.1).

The constant term of $N_{a,b,c}(x)$ is $1$. Consequently, a palindromic
polynomial must be monic. Ding and Dong \cite[Lemma~2.3(a)]{ding2019} classify the monic
parameters. In the sorted range $a\le b\le c$, their condition is equivalent
to

\nopagebreak[4]
\[
c=a+b-1-2d,
\qquad
0\le d\le\left\lfloor\frac{a-1}{2}\right\rfloor.
\tag{2.2}
\]

For these parameters, put

\nopagebreak[4]
\[
H=a+b-2
\]

and

\nopagebreak[4]
\[
B_d=\{(i,j)\in Q:d\le i+j\le H-d\}.
\tag{2.3}
\]

The unique maximum support is $B_d$, and its size is

\nopagebreak[4]
\[
w=|B_d|=ab-d(d+1).
\tag{2.4}
\]

Thus non-monic parameters are already excluded. We shall compare
$[x]N_{a,b,c}(x)$ and $[x^{w-1}]N_{a,b,c}(x)$ throughout the monic range.
The first coefficient is immediate:

\nopagebreak[4]
\[
[x]N_{a,b,c}(x)=abc.
\tag{2.5}
\]

\section{Path compensation}


Assume in this section that

\nopagebreak[4]
\[
1\le d\le\left\lfloor\frac{a-1}{2}\right\rfloor.
\]

Let

\nopagebreak[4]
\[
T_0=Q\setminus B_d.
\]

Every full northeast path contains exactly $2d$ points of $T_0$. Let $S$ be
a binary support of size $w-1$, and write its complement uniquely as

\nopagebreak[4]
\[
Q\setminus S=(T_0\setminus R)\cup A,
\tag{3.1}
\]

where

\nopagebreak[4]
\[
R\subseteq T_0,\qquad A\subseteq B_d,\qquad |A|=|R|+1.
\tag{3.2}
\]

The binary matrix $\mathbf1_S$ satisfies $L(\mathbf1_S)\le c$ if and only if

\nopagebreak[4]
\[
|P\cap A|\ge |P\cap R|
\quad\text{for every full northeast path }P.
\tag{3.3}
\]

We call $(R,A)$ a path-compensation pair.

For $0\le r\le H$, define the rank diagonal

\nopagebreak[4]
\[
D_r=\{(i,j)\in Q:i+j=r\}
\]

and, for $0\le r\le d$, enumerate it as

\nopagebreak[4]
\[
D_r=\{x_{r,0},\ldots,x_{r,r}\}.
\]

Let

\nopagebreak[4]
\[
R^-=R\cap(D_0\cup\cdots\cup D_{d-1}).
\]

For $x_{r,k}$ define

\nopagebreak[4]
\[
v_r(k)=
\max_P |P\cap R^-|,
\tag{3.4}
\]

where $P$ ranges over paths from $(0,0)$ to $x_{r,k}$. Then

\nopagebreak[4]
\[
v_r(k)=\mathbf1_{x_{r,k}\in R^-}
+\max\{v_{r-1}(k-1),v_{r-1}(k)\},
\tag{3.5}
\]

with nonexistent predecessors omitted. Define

\nopagebreak[4]
\[
G^-(R^-)=\sum_{k=0}^{d}v_d(k).
\tag{3.6}
\]

\Needspace{0.25\textheight}
\begin{lemma}

If $r_0$ is the lowest nonempty rank of $R^-$, then

\nopagebreak[4]
\[
G^-(R^-)\ge |R^-|+(d-r_0).
\tag{3.7}
\]

Moreover,

\nopagebreak[4]
\[
G^-(R^-)=|R^-|+1
\tag{3.8}
\]

if and only if $R^-$ is a nonempty consecutive interval in $D_{d-1}$.

\end{lemma}

\begin{proof}

For a nonnegative vector $u=(u_0,\ldots,u_m)$, set

\nopagebreak[4]
\[
E(u)=
(u_0,\max(u_0,u_1),\ldots,\max(u_{m-1},u_m),u_m).
\]

Then

\nopagebreak[4]
\[
\sum E(u)-\sum u
=u_m+\sum_{j=0}^{m-1}(u_j-u_{j+1})_+.
\tag{3.9}
\]

Starting at the lowest occupied rank and applying (3.5)--(3.9) at every
subsequent rank proves (3.7). If $r_0<d-1$, at least two strict frontier
expansions occur. If the lowest occupied rank has more than one consecutive
component, the first expansion already creates at least two new frontier
units. Any additional removal on a higher outer rank creates another unit in
(3.5). Hence (3.8) forces a single interval in $D_{d-1}$. Conversely, an
interval of length $\ell$ in $D_{d-1}$ produces an interval of length
$\ell+1$ on $D_d$, so equality holds.
\end{proof}


Rotating the rectangle through $180$ degrees gives the analogous high-end
quantities $R^+$ and $G^+(R^+)$.

\Needspace{0.25\textheight}
\begin{lemma}

Every path-compensation pair satisfies

\nopagebreak[4]
\[
|A|\ge G^-(R^-)+G^+(R^+).
\tag{3.10}
\]

Consequently, $R^-$ and $R^+$ cannot both be nonempty.

\end{lemma}

\begin{proof}

Write

\nopagebreak[4]
\[
y_p=(p,d-p),\qquad
z_p=(a-1-d+p,b-1-p),
\qquad 0\le p\le d.
\]

The inequality $2d\le a-1\le b-1$ gives $y_p\le z_p$. Use one fixed step
word between each pair $(y_p,z_p)$. Its $d+1$ translates form pairwise
vertex-disjoint central channels.

At $y_p$, attach a low-end prefix attaining $v_d(p)$; at $z_p$, attach a
high-end suffix attaining the corresponding rotated maximum. Applying
(3.3) to these $d+1$ full paths and summing, every point of $A\subseteq B_d$
is counted at most once, while the right side is exactly
$G^-(R^-)+G^+(R^+)$. This proves (3.10).

If both ends contain removals, Lemma 3.1 and its rotation give

\nopagebreak[4]
\[
|A|\ge |R^-|+|R^+|+2=|R|+2,
\]

contrary to (3.2).
\end{proof}


\section{Minimum separators in the central band}


We next classify the minimum vertex separators used to compensate a single
end interval.

\Needspace{0.25\textheight}
\begin{lemma}

Let $U\subseteq D_{d+1}$ with $|U|=q\le d+1$. There are $q$ pairwise
vertex-disjoint northeast paths joining the points of $U$ to $q$ distinct
points of $D_{H-d}$.

\end{lemma}

\begin{proof}

Proceed one rank at a time. Before the final rank, every intermediate
diagonal in the central band has at least $d+2$ points. For a proper subset
$Y$ of any such diagonal, its complete upper shadow has size at least
$|Y|$. The only way a shadow can lose one point is for $Y$ to fill an entire
diagonal in the shrinking part of the rectangle. This cannot happen because
there are only $q\le d+1$ active points.

Thus the bipartite graph between the current points and their successors on
the next diagonal satisfies Hall's condition. Choose a matching at each
rank. Composing the matchings gives the required disjoint paths.
\end{proof}


\Needspace{0.25\textheight}
\begin{corollary}

Let $\varnothing\ne X\subseteq D_d$ with $|X|=k\le d$. Any vertex set
disjoint from $X$ that meets every path from $X$ to $D_{H-d}$ has at least
$k+1$ points.

\end{corollary}

\begin{proof}

The complete upper shadow of $X$ satisfies

\nopagebreak[4]
\[
|\partial^+X|
=|X|+\#\{\text{consecutive components of }X\}
\ge k+1.
\]

Choose $k+1$ shadow points and apply Lemma 4.1. Extend the resulting paths
down by one edge to $X$. They can share starting points in $X$, but outside
$X$ they are pairwise vertex-disjoint. A separator disjoint from $X$ must
therefore use at least one distinct point on each path.
\end{proof}


\Needspace{0.25\textheight}
\begin{lemma}

Let $J\subseteq D_d$ be a nonempty consecutive interval of size $m$, and
let $A\subseteq B_d$ meet every path from $J$ to $D_{H-d}$. If $|A|=m$,
then

\nopagebreak[4]
\[
A=J,
\tag{4.1}
\]

except that when $J=D_d$ one may also have

\nopagebreak[4]
\[
A=D_{H-d}.
\tag{4.2}
\]

\end{lemma}

\begin{proof}

Put $X=J\setminus A$. If $X$ is nonempty and $|X|=k\le d$, then $A\cap J$
uses $m-k$ points, while Corollary 4.2 forces at least $k+1$ further points.
This exceeds the budget $m$. Therefore, if $m<d+1$, one must have $A=J$.

Now let $m=d+1$, so $J=D_d$. The preceding argument shows that $A$ either
equals $D_d$ or avoids $D_d$ completely. If $A$ uses a nonempty proper
subset of $D_{H-d}$, apply the rotated Corollary 4.2 to
$D_{H-d}\setminus A$; the same budget contradiction results. Thus $A$ is
the whole upper boundary or avoids both boundaries.

In the latter case, apply the rank-by-rank Hall argument to the full
$D_{d+1}$ and $D_{H-d-1}$. It gives $d+2$ internally disjoint paths. After
extending them to the two boundary diagonals, they can share vertices only
on boundaries that $A$ avoids. Hence an internal separator needs at least
$d+2$ points, again a contradiction.
\end{proof}


\section{Binary supports one below maximum}


\Needspace{0.25\textheight}
\begin{theorem}

Suppose $d\ge1$. The binary supports $S\subseteq Q$ satisfying

\nopagebreak[4]
\[
|S|=w-1,\qquad L(\mathbf1_S)\le c
\]

are exactly the following:

\begin{enumerate}
\item $S=B_d\setminus\{z\}$ for $z\in B_d$;
\item at either end, let $I$ be a nonempty consecutive interval in the outer
diagonal of $T_0$, and let $J(I)$ be its complete shadow on the adjacent
inner diagonal; then

\nopagebreak[4]
\[
S=Q\setminus\bigl((T_0\setminus I)\cup J(I)\bigr);
\]

\item the two cross-end full-diagonal exchanges

\nopagebreak[4]
\[
S=Q\setminus\bigl((T_0\setminus D_{d-1})\cup D_{H-d}\bigr)
\]

and

\nopagebreak[4]
\[
S=Q\setminus\bigl((T_0\setminus D_{H-d+1})\cup D_d\bigr).
\]

\end{enumerate}

Their total number is

\nopagebreak[4]
\[
ab+2.
\tag{5.1}
\]

\end{theorem}

\begin{proof}

Use the complement representation (3.1). If $R=\varnothing$, then
$|A|=1$, giving the first family.

Assume $R\ne\varnothing$. Lemma 3.2 puts all removals at one end.
Lemma 3.1 and $|A|=|R|+1$ force $R$ to be a single consecutive interval
$I$ in the outermost diagonal $D_{d-1}$, with no other removals.

Let

\nopagebreak[4]
\[
J=\partial^+I\subseteq D_d.
\]

Every path from $J$ to $D_{H-d}$ can be extended downward through a point
of $I$ and then to the source. Condition (3.3) therefore says that $A$ is a
separator for all such paths. Moreover,

\nopagebreak[4]
\[
|A|=|I|+1=|J|.
\]

Lemma 4.3 gives $A=J$, except that when $I=D_{d-1}$ it also permits
$A=D_{H-d}$. These are the same-end exchanges and one cross-end exchange.
Rotation gives the high-end versions and the other cross-end exchange.
Every listed pair directly satisfies (3.3).

There are $w$ supports in the first family. At each end, the number of
nonempty intervals in a diagonal of length $d$ is

\nopagebreak[4]
\[
\sum_{\ell=1}^{d}(d-\ell+1)=\frac{d(d+1)}2.
\]

Adding both ends and the two cross-end exchanges gives

\nopagebreak[4]
\[
w+d(d+1)+2=ab+2.
\]

This proves (5.1).
\end{proof}


\section{Weight increments and the coefficient for \texorpdfstring{$d\ge1$}{d >= 1}}


A matrix counted by $[x^{w-1}]N_{a,b,c}(x)$ may have entries larger than
one. We now classify all such possibilities.

First consider

\nopagebreak[4]
\[
S=B_d\setminus\{z\}.
\]

Increasing the entry at $u\in S$ from $1$ to $2$ is legal if and only if
every saturated central-band path through $u$ also passes through $z$.
Equivalently, $z$ dominates every source-to-$u$ path or every
$u$-to-sink path in the directed grid induced by $B_d$.

\Needspace{0.25\textheight}
\begin{lemma}

The legal ordered pairs $(z,u)$ are exactly those on one of the four exposed
boundary chains, with $z$ separating the corresponding boundary endpoint
from $u$. Their number is

\nopagebreak[4]
\[
(a-d)(a-d-1)+(b-d)(b-d-1).
\tag{6.1}
\]

No matrix with two incremented points or an increment of at least two is
legal.

\end{lemma}

\begin{proof}

If $u\in D_d$, there is no central-band point of smaller rank that can
dominate it. If $\operatorname{rank}(u)>d$ and $u$ is not on a source-side
exposed chain, two adjacent points of $D_d$ reach $u$ along two paths that
meet only at $u$. Hence no $z\ne u$ lies on every source-to-$u$ path. On an
exposed chain, the path from its unique reachable boundary endpoint to $u$
is forced, and its intermediate points are precisely the dominators.
Rotation gives the sink-side statement, including the top boundary rank.

The four chain lengths are

\nopagebreak[4]
\[
a-d,\quad a-d,\quad b-d,\quad b-d.
\]

Counting separated ordered pairs gives

\nopagebreak[4]
\[
2\binom{a-d}{2}+2\binom{b-d}{2},
\]

which is (6.1).

For fixed $z$, every legal $u$ lies on one exposed chain. Any two such points
occur together on a path that, after deleting $z$, has only $c-1$ basic
support points. Two unit increments, or two units at one point, raise its
weight to $c+1$. Thus no further increment is possible.
\end{proof}


\Needspace{0.25\textheight}
\begin{lemma}

Every point of every exchange support in Theorem 5.1 lies on a path with
exactly $c$ support points. Hence exchange supports admit no positive weight
increment.

\end{lemma}

\begin{proof}

Consider a low-end interval exchange. Let

\nopagebreak[4]
\[
I\subseteq D_{d-1},
\qquad
J=\partial^+I\subseteq D_d.
\]

If $u\in I$, extend a prefix through $u$ to $J$ and then to the sink; the
path meets both $I$ and $J$ once. If $u\in D_d\setminus J$, enter $u$ from
a predecessor in $D_{d-1}\setminus I$; the path avoids both sets.

Finally, suppose $\operatorname{rank}(u)>d$. The points of $D_d$ that can
reach $u$ form an interval. If this interval meets $J$, enter through an
intersection point and one of its predecessors in $I$. If it does not meet
$J$, enter through a point outside $J$ and a predecessor outside $I$. The
resulting path still passes through $u$ and meets $I,J$ equally.

Thus every support point lies on a path satisfying

\nopagebreak[4]
\[
|P\cap I|=|P\cap J|.
\]

Such a path has exactly $c$ support points after the exchange. The high-end
case follows by rotation. For a cross-end full exchange, every full path
meets the removed and added diagonals once, so every path is already
saturated.
\end{proof}


Combining Theorem 5.1 and Lemmas 6.1--6.2 yields the exact coefficient.

\Needspace{0.25\textheight}
\begin{theorem}

For

\nopagebreak[4]
\[
3\le a\le b,\qquad
1\le d\le\left\lfloor\frac{a-1}{2}\right\rfloor,
\]

put

\nopagebreak[4]
\[
c=a+b-1-2d,\qquad w=ab-d(d+1).
\]

Then

\nopagebreak[4]
\[
\boxed{
[x^{w-1}]N_{a,b,c}(x)
=ab+2+(a-d)(a-d-1)+(b-d)(b-d-1).
}
\tag{6.2}
\]

Moreover,

\nopagebreak[4]
\[
[x^{w-1}]N_{a,b,c}(x)<abc.
\tag{6.3}
\]

\end{theorem}

\begin{proof}

Theorem 5.1 counts $ab+2$ binary supports. Lemma 6.1 contributes the two
families of single increments in (6.1), and Lemma 6.2 excludes increments
on every other support. This proves (6.2).

Expanding the difference gives

\nopagebreak[4]
\[
\begin{aligned}
abc-[x^{w-1}]N
={}&(a-1)b^2+(a-1)(a-1-2d)b\\
&-(a-d)^2-d^2+a-2d-2.
\end{aligned}
\tag{6.4}
\]

The second term is nonnegative. Since

\nopagebreak[4]
\[
(a-d)^2+d^2\le a^2,\qquad
2d\le a-1,\qquad b\ge a,
\]

we obtain

\nopagebreak[4]
\[
abc-[x^{w-1}]N
\ge(a-2)a^2-1>0.
\]

This proves (6.3).
\end{proof}


\section{The zero-depth case}


It remains to treat the monic parameter

\nopagebreak[4]
\[
d=0,\qquad c=a+b-1,\qquad w=ab.
\]

A support of size $ab-1$ is $Q\setminus\{z\}$. Put weight $1$ on this
support and denote all additional weights by $E$. For every full path $P$,
the constraint is

\nopagebreak[4]
\[
\sum_{u\in P}E_u\le\mathbf1_{z\in P}.
\tag{7.1}
\]

If $z$ is the source or sink, every path contains $z$. The nonzero support
of $E$ may then be any nonempty antichain other than $\{z\}$. The number of
antichains in an $a$ by $b$ rectangle is

\nopagebreak[4]
\[
\binom{a+b}{a}.
\]

Thus each endpoint contributes

\nopagebreak[4]
\[
\binom{a+b}{a}-2
\]

additional matrices beyond the binary matrix.

For non-endpoint $z$, an increment is possible only on the same exposed
outer boundary, with $z$ separating the endpoint from the incremented
point. As in Lemma 6.1, no two increments and no larger increment are
possible. The four boundary contributions total

\nopagebreak[4]
\[
2\binom{a-1}{2}+2\binom{b-1}{2}.
\]

Including the $ab$ binary supports gives:

\Needspace{0.25\textheight}
\begin{proposition}

For $3\le a\le b$,

\nopagebreak[4]
\[
\begin{aligned}
[x^{ab-1}]N_{a,b,a+b-1}(x)
={}&ab+2\left(\binom{a+b}{a}-2\right)\\
&+(a-1)(a-2)+(b-1)(b-2)\\
={}&2\binom{a+b}{a}+a^2+b^2+ab-3a-3b.
\end{aligned}
\tag{7.2}
\]

This coefficient satisfies

\nopagebreak[4]
\[
[x^{ab-1}]N_{a,b,a+b-1}(x)>ab(a+b-1).
\tag{7.3}
\]

\end{proposition}

\begin{proof}

Only (7.3) remains. Its left side minus right side is

\nopagebreak[4]
\[
2\binom{a+b}{a}-(a+b)(ab-a-b+3).
\tag{7.4}
\]

Set

\nopagebreak[4]
\[
R_{a,b}=
\frac{2\binom{a+b}{a}}
{(a+b)((a-1)(b-1)+2)}.
\]

For fixed $a\ge3$, comparison of consecutive ratios gives

\nopagebreak[4]
\[
\begin{aligned}
&\frac{\binom{a+b+1}{a}}{\binom{a+b}{a}}
-\frac{(a+b+1)((a-1)b+2)}
{(a+b)((a-1)(b-1)+2)}\\
&\qquad=
\frac{(a-2)(a-1)(b-1)(a+b+1)}
{(a+b)(b+1)(ab-a-b+3)}>0.
\end{aligned}
\]

Thus $R_{a,b}$ increases with $b$. On the diagonal, let

\nopagebreak[4]
\[
L_a=2\binom{2a}{a},
\qquad
U_a=2a((a-1)^2+2).
\]

The base case is $L_3=40>36=U_3$, and

\nopagebreak[4]
\[
\frac{L_{a+1}}{L_a}-\frac{U_{a+1}}{U_a}
=
\frac{(a-1)^2(3a^2-2a-2)}
{a(a+1)(a^2-2a+3)}>0.
\]

Hence $R_{a,b}>1$ for all $3\le a\le b$, which is precisely (7.3).
\end{proof}


\section{Proof of the classification}


\begin{proof}[Proof of Theorem 1.2]

If $N_{a,b,c}(x)$ is not monic, it cannot be palindromic. In the monic
range, write $c=a+b-1-2d$ as in (2.2).

If $d=0$, Proposition 7.1 gives

\nopagebreak[4]
\[
[x^{w-1}]N_{a,b,c}(x)>[x]N_{a,b,c}(x).
\]

If $d\ge1$, Theorem 6.3 gives the strict reverse inequality. In either case,
the coefficients paired by palindromicity are unequal. Therefore
$N_{a,b,c}(x)$ is not palindromic.
\end{proof}


\begin{proof}[Proof of Theorem 1.1]

When $a=1$, the number of $k$-element antichains in
$[b]\times[c]$ is

\nopagebreak[4]
\[
\binom{b}{k}\binom{c}{k}.
\]

After sorting $b\le c$, its top coefficient is $\binom{c}{b}$, so the
polynomial is monic, and hence possibly palindromic, only when $b=c$.
For $b=c=r$, the coefficients are $\binom{r}{k}^2$ and are palindromic.

When $a=2$, the complete classification in \cite{carptopus2026} states that the polynomial is
palindromic exactly when $c=b+1$. When $a\ge3$, Theorem 1.2 excludes every
parameter. This proves the stated list.
\end{proof}


\Needspace{0.25\textheight}
\begin{corollary}

Ding--Dong Conjecture B holds for the complete rectangular minuscule family.

\end{corollary}

\begin{proof}

Theorem 1.1 lists all palindromic cases. For the one-layer square, Ding and
Dong \cite[Proposition~2.4]{ding2019} give the gamma expansion

\nopagebreak[4]
\[
\sum_{k=0}^{r}\binom{r}{k}^2x^k
=
\sum_{j=0}^{\lfloor r/2\rfloor}
\frac{r!}{j!\,j!\,(r-2j)!}\,
x^j(1+x)^{r-2j}.
\]

Every allowed gamma coefficient is strictly positive. Reference \cite{carptopus2026} proves
strict gamma-positivity for every palindromic two-layer case. There are no
palindromic higher-layer cases.
\end{proof}


\section{Scope and verification}


The proof is an all-parameter argument. Finite computation was used only for
calibration and adversarial testing:

\begin{enumerate}
\item the Ding--Dong transfer recurrence reproduces their published example and
scans small monic parameters;
\item direct complement enumeration verifies the binary support count and every
possible single-cell unit increment for several $d=1,2$ cases, and retains
a $d=3$ negative control that disproved an earlier incomplete formula;
\item an independent row-frontier dynamic program verifies the binary support
count in four $d=3$ cases and obtains the value $75$ for $(7,7,7)$ under
the entry bound $0,1,2$.

\end{enumerate}

These checks do not replace Lemmas 3.1--6.2. In particular, solver timeouts or
finite absence of counterexamples are not used as existence or nonexistence
proofs. Entries at least $3$ and simultaneous increments are excluded by
Lemmas 6.1--6.2, not by the finite scripts.

The result settles only the rectangular minuscule branch of Conjecture B. It
does not claim a classification for the other connected minuscule posets.

\section*{Acknowledgement of AI assistance}


The author used Codex as an AI-assisted research and writing tool for
literature organization, symbolic and finite verification, adversarial proof
audits, and drafting. Carptopus is the sole author and assumes responsibility
for every statement and proof.

\begin{thebibliography}{9}
\footnotesize
\raggedright

\bibitem{ding2019}
J.~Ding and C.-P.~Dong,
\emph{Antichain generating polynomials of posets},
arXiv:1905.06692v1 (2019).
\href{https://arxiv.org/abs/1905.06692}{arXiv:1905.06692}.

\bibitem{amanov2023}
A.~Amanov and D.~Yeliussizov,
\emph{MacMahon's statistics on higher-dimensional partitions},
Forum Math. Sigma 11 (2023), e63.
\href{https://doi.org/10.1017/fms.2023.61}{doi:10.1017/fms.2023.61}.

\bibitem{carptopus2026}
Carptopus,
\emph{Real-rootedness, palindromicity, and gamma-positivity of antichain
polynomials for $[2]\times[m]\times[n]$},
v0.2-beta (2026).
\href{https://doi.org/10.5281/zenodo.22112708}{doi:10.5281/zenodo.22112708}.

\end{thebibliography}

\end{document}
